\subsection{Another interpretation of the topology of the space of closed subgroups}

Let F be the set of closed subsets of G; one defines a Hausdorff uniform structure on F in the following way: for every compact subset K of G and every neighborhood V of e in G, let \(\mathrm{P}(\mathrm{K},\mathrm{V})\) be the set of pairs \((\mathrm{X},\mathrm{Y})\) of elements of F such that both

\[
\mathrm{X} \cap \mathrm{K} \subset \mathrm{VY} \quad \text { and } \quad \mathrm{Y} \cap \mathrm{K} \subset \mathrm{VX}.\tag{9}
\]

Let us show that the set of \(\mathrm{P}(\mathbf{K},\mathbf{V})\) is a fundamental system of entourages for a Hausdorff uniform structure \(\mathcal{U}\) on \(\mathfrak{F}\) . The axioms \((\mathrm{U}_{\mathrm{I}}^{\prime})\) and \((\mathrm{U}_{\mathrm{II}}^{\prime})\) of GT, II, §1, No.~1 are obviously satisfied; moreover, the relations \(\mathbf{K} \subset \mathbf{K}'\) and \(\mathbf{V}' \subset \mathbf{V}\) imply \(\mathrm{P}(\mathbf{K}',\mathbf{V}') \subset \mathrm{P}(\mathbf{K},\mathbf{V})\) ; to verify \((\mathrm{U}_{\mathrm{III}}^{\prime})\) , one can therefore limit oneself to the case that \(\mathbf{V}\) is a symmetric compact neighborhood of \(e\) , so that VK is compact. Suppose that \((\mathbf{X},\mathbf{Y}) \in \mathrm{P}(\mathrm{VK},\mathbf{V})\) and \((\mathbf{Y},\mathbf{Z}) \in \mathrm{P}(\mathrm{VK},\mathbf{V})\) ; then \(\mathbf{X} \cap \mathbf{K} \subset \mathbf{X} \cap \mathrm{VK} \subset \mathrm{VY}\) , and if \(y \in \mathbf{Y}\) is such that \(vy \in \mathbf{K}\) for some \(v \in \mathbf{V}\) , then necessarily \(y \in \mathrm{VK}\) , therefore

\[
\mathrm{X} \cap \mathrm{K} \subset \mathrm{V} (\mathrm{Y} \cap \mathrm{VK});
\]

on the other hand, \(Y \cap VK \subset VZ\) , whence \(X \cap K \subset V^{2}Z\) , and one shows similarly that \(Z \cap K \subset V^{2}X\) , which proves \((U_{III}^{\prime})\) . Finally, if X,Y are two distinct elements of F, there exists for example a point \(a \in X\) such that \(a \notin Y\) , hence a symmetric compact neighborhood V of e such that \(Va \cap Y = \varnothing\) , that is, \(a \notin VY\) ; a fortiori \((X,Y) \notin P(Va,V)\) , which completes the proof of our assertion.

This established, let us consider on the set \(\Sigma\) of closed subgroups of \(G\) the topology \(\mathcal{T}\) induced by the topology of the uniform space \(\mathfrak{F}\) just defined. We shall see that this topology is identical to the topology defined in No.~3. It will suffice to prove that the mapping \(\alpha \mapsto H_{\alpha}\) of \(\Gamma\) into \(\Sigma\) is continuous when \(\Sigma\) is equipped with the topology \(\mathcal{T}\) : for, the same will then be true of the restriction of this mapping to \(\Gamma_{\varphi}\) (with notations as in No.~1, Prop. 2), which is bijective; but since \(\Gamma_{\varphi}\) is compact and the topology \(\mathcal{T}\) is separated, the mapping \(\alpha \mapsto H_{\alpha}\) of \(\Gamma_{\varphi}\) into \(\Sigma\) will then be a homeomorphism.

Thus let \(\alpha_0\) be a point of \(\Gamma\) and let \(\Phi\) be a filter on \(\Gamma\) that converges to \(\alpha_0\) ; we are to show that, with respect to \(\Phi\) , \(\mathrm{H}_{\alpha}\) tends to \(\mathrm{H}_{\alpha_0}\) for the topology \(\mathcal{T}\) . Let \(\mathbf{K}\) be a compact subset of \(\mathbf{G}\) , \(\mathbf{V}\) a symmetric compact neighborhood of \(e\) in \(\mathbf{G}\) ; for every \(x \in \mathrm{H}_{\alpha_0} \cap \mathrm{K}\) , there exists a set \(\mathrm{M}(x) \in \Phi\) such that for every \(\alpha \in \mathrm{M}(x)\) , one has \(\mathrm{V}x \cap \mathrm{H}_{\alpha} \neq \emptyset\) (No.~1, Lemma 2), whence \(\mathrm{V}x \subset \mathrm{V}^2\mathrm{H}_{\alpha}\) ; on covering \(\mathrm{H}_{\alpha_0} \cap \mathrm{K}\) by a finite number of sets \(\mathrm{V}x_i\) , one sees that if \(\mathrm{M} = \bigcap \mathrm{M}(x_i)\) , then \(\mathrm{H}_{\alpha_0} \cap \mathrm{K} \subset \mathrm{V}^2\mathrm{H}_{\alpha}\) for every \(\alpha \in \mathrm{M}\) .

Conversely, suppose that there existed an open neighborhood U of e in G such that, for every set \(L \in \Phi\) , there is at least one \(\alpha \in L\) for which \(H_{\alpha} \cap K \not\subset UH_{\alpha_{0}}\) ; if \(\omega(L)\) is the set of \(\alpha \in L\) having this property, the \(\omega(L)\) would form a base of a filter \(\Phi'\) on \(\Gamma\) finer than \(\Phi\) , and, for every \(\alpha\) belonging to the union E of the \(\omega(L)\) for \(L \in \Phi\) , there would exist a \(t_{\alpha} \in H_{\alpha} \cap K\) not belonging to \(UH_{\alpha_0}\) ; for \(\alpha \notin E\) , take for \(t_{\alpha}\) any point of \(H_{\alpha}\) . Since \(K \cap \mathbb{C}(UH_{\alpha_0})\) is compact, there would exist a cluster point \(s\) of \(\alpha \mapsto t_{\alpha}\) with respect to \(\Phi'\) , belonging to \(K \cap \mathbb{C}(UH_{\alpha_0})\) ; but since \(\Phi'\) converges to \(\alpha_0\) in \(\Gamma\) , this contradicts Lemma 3 of No.~1.

